\chapter{Attention Mechanisms}
\label{ch:09}

% ─── Learning Goals ───
\begin{keyinsight}
After reading this chapter, you will be able to:
\begin{enumerate}
  \item Understand the mathematical motivation for the attention mechanism as a differentiable dictionary lookup.
  \item Derive the scaled dot-product attention equation and explain the need for the scaling factor $\frac{1}{\sqrt{d_k}}$.
  \item Construct and apply a causal mask for autoregressive sequence modeling.
  \item Explain how Multi-Head Attention allows a model to represent multiple, independent relationships between tokens.
\end{enumerate}
\end{keyinsight}

% ─── Big Picture ───
\section{Big Picture}
\label{ch:09:sec:big_picture}

As we saw in Chapter 8, older architectures like RNNs compress the entire history of a sequence into a single, fixed-size vector. This creates an unmanageable information bottleneck as sequences grow longer.

The \textbf{Attention Mechanism} solves this by discarding compression entirely. Instead of trying to remember the past, an attention layer looks at the \emph{exact} past tokens at every step. It acts like a powerful, differentiable search engine: each token generates a "query" about what information it needs, scans all previous tokens' "keys" to find matches, and then extracts their "values". 

This mechanism is the core of the Transformer architecture. Every other component in an LLM exists solely to feed, normalize, or process the outputs generated by attention.

% ─── Formal Setup ───
\section{Formal Setup and Notation}
\label{ch:09:sec:setup}

Let $\mathbf{X} \in \R^{T \times d_m}$ be a sequence of $T$ tokens, where $d_m$ is the model dimension. 
The attention mechanism uses three projection matrices to map the input $\mathbf{X}$ into three distinct spaces:
\begin{itemize}
    \item \textbf{Queries} ($\vQ = \mathbf{X} \vW_Q \in \R^{T \times d_k}$): What each token is looking for.
    \item \textbf{Keys} ($\vK = \mathbf{X} \vW_K \in \R^{T \times d_k}$): What each token contains.
    \item \textbf{Values} ($\vV = \mathbf{X} \vW_V \in \R^{T \times d_v}$): The actual information to be extracted.
\end{itemize}
Typically, $d_k = d_v$. We will assume $d_k$ is the attention dimension.

% ─── Main Ideas ───
\section{Scaled Dot-Product Attention}
\label{ch:09:sec:main}

The core operation computes the interaction between every query and every key using a dot product, applies a softmax to convert these unbounded scores into probabilities, and uses them to compute a weighted sum of the values.

The canonical equation is:
\begin{equation}
\label{eq:09:attention}
\Attn(\vQ, \vK, \vV) = \softmax\left(\frac{\vQ \vK^\T}{\sqrt{d_k}}\right) \vV
\end{equation}

\subsection{The Pre-Softmax Scaling Factor}
Why divide by $\sqrt{d_k}$? Assume the elements of $\vQ$ and $\vK$ are independent random variables with zero mean and variance $1$. The dot product of a query vector $\mathbf{q} \in \R^{d_k}$ and a key vector $\mathbf{k} \in \R^{d_k}$ is $\mathbf{q} \cdot \mathbf{k} = \sum_{i=1}^{d_k} q_i k_i$.
Because this is a sum of $d_k$ independent terms, its variance is $d_k$. 

If $d_k = 64$, the variance of the dot product is $64$, meaning the pre-softmax scores will have huge magnitudes. The softmax function $\frac{\exp(x_i)}{\sum \exp(x_j)}$ will be pushed deep into regions where the gradient $\approx 0$, stalling training. Dividing the dot product by $\sqrt{d_k}$ forces the variance back to $1$.

\section{Causal Masking}

Equation \cref{eq:09:attention} computes the interaction between \emph{all} tokens. However, in an autoregressive language model (Chapter~6), predicting token $x_t$ must only depend on $x_{<t}$. If token $t$ attends to token $t+1$, the model is ``cheating'' by looking into the future during training.

We prevent this by adding a \textbf{mask matrix} $\mathbf{M} \in \R^{T \times T}$ before the softmax:
\begin{equation}
M_{ij} = \begin{cases} 
0 & \text{if } i \geq j \\
-\infty & \text{if } i < j 
\end{cases}
\end{equation}

Thus, the causal attention equation becomes:
\begin{equation}
\label{eq:09:causal_attn}
\Attn(\vQ, \vK, \vV) = \softmax\left(\frac{\vQ \vK^\T}{\sqrt{d_k}} + \mathbf{M}\right) \vV
\end{equation}
Because $\exp(-\infty) = 0$, the attention weights for future tokens are strictly zero.

\section{Multi-Head Attention (MHA)}

A single attention mechanism computing a single weighted sum of values is restrictive. What if a "verb" token needs to attend to its "subject" token to check plural agreement, but simultaneously needs to attend to an "adverb" token to check semantics? 

We solve this by running $h$ independent attention mechanisms (heads) in parallel.
\begin{enumerate}
    \item Project $\mathbf{X}$ into $h$ different sets of queries, keys, and values.
    \item Compute Equation~\cref{eq:09:causal_attn} for each head independently.
    \item Concatenate the $h$ output matrices.
    \item Linearly project the concatenated matrix back to the model dimension $d_m$ using a matrix $\vW_O$.
\end{enumerate}

To keep parameters and compute constant regardless of $h$, we split the dimensions: $d_k = d_v = \frac{d_m}{h}$. If $d_m = 4096$ and $h = 32$, each head operates in a 128-dimensional subspace.

% ─── Worked Example ───
\section{Worked Example: Attention Matrix Shapes}
\label{ch:09:sec:example}

Given: batch size $B=4$, sequence length $T=10$, model dimension $d_m = 512$, and heads $h=8$.
What are the shapes of all intermediate tensors when executed in PyTorch?

1. The input $\mathbf{X}$ has shape $(B, T, d_m) = (4, 10, 512)$.
2. We project to $\vQ, \vK, \vV$. The query projection matrix $\vW_Q$ has shape $(512, 512)$. (Because $h \times d_k = 8 \times 64 = 512$).
3. The resulting $\vQ$ has shape $(4, 10, 512)$. We reshape and transpose it to isolate the heads: $(B, h, T, d_k) = (4, 8, 10, 64)$.
4. We compute $\vQ \vK^\T$. We want the dot product over the $d_k$ dimension.
   Shape: $(4, 8, 10, 64) \times (4, 8, 64, 10) \rightarrow (4, 8, 10, 10)$. 
   This is the $10 \times 10$ attention score matrix for each head and batch.
5. Apply causal mask (shape $10 \times 10$ broadcasted) and softmax.
6. Multiply by $\vV$: $(4, 8, 10, 10) \times (4, 8, 10, 64) \rightarrow (4, 8, 10, 64)$.
7. Reshape (concatenate heads): $(4, 10, 512)$.
8. Final projection by $\vW_O$ (shape $512 \times 512$) yields $(4, 10, 512)$.

Notice that the input and output shapes of the entire Multi-Head Attention block are identical.

% ─── Systems Consequences ───
\section{Systems and Implementation Consequences}
\label{ch:09:sec:systems}

The attention matrix $\mathbf{S} = \vQ \vK^\T$ has shape $T \times T$. This establishes the fundamental scaling law of the vanilla Transformer: memory and compute scale \textbf{quadratically} $O(T^2)$ with the sequence length. 

For $T=1000$, $1000 \times 1000$ requires minimal memory. For a document where $T=100,000$, the $T \times T$ matrix alone would swallow gigabytes of VRAM per batch per layer. As we will see in Chapter 16, avoiding the materialization of this $T \times T$ matrix using IO-aware kernels (FlashAttention) is the most critical systems optimization in modern ML.

% ─── Neuroscience Analogy ───
\begin{analogybox}{The Cocktail Party Effect and Selective Attention}
  \paragraph{Where the analogy helps.}
  At a loud party, your brain can dynamically suppress background noise and isolate one specific person's voice (the "target"). The Transformer's attention mechanism behaves similarly: given a word (the query), the softmax distribution suppresses irrelevant context words (zeroing out their weights) and attends sharply to the geometrically relevant words (the keys), routing only their information forward.

  \paragraph{Where the analogy breaks.}
  Human selective attention happens in real-time sequential processing and heavily involves physical sensory gating via the thalamus. Neural network attention operates instantly across all past tokens simultaneously via dense linear algebra. Furthermore, human attention is tightly bottlenecked (we can generally only focus on one conversational stream), whereas Multi-Head Attention computes dozens of independent attentional streams in parallel without interference.

  \paragraph{What intuition to keep.}
  Attention is a routing mechanism. Instead of pushing all text through rigid, fixed pathways, the data itself determines what connections matter for each specific word.
\end{analogybox}


% ─── Misconceptions ───
\section{Common Misconceptions}
\label{ch:09:sec:misconceptions}

\begin{enumerate}
  \item \textbf{``Attention is all you need.''} The title of the paper is catchy, but false. Attention merely routes information. If you remove the heavy Feed-Forward Networks (MLPs) from the Transformer layer, the model entirely loses its ability to store facts and perform complex non-linear reasoning. 
  \item \textbf{``Attention weights show exactly how the model thinks.''} Visualizing the softmax weights $10 \times 10$ matrix gives a nice graphic associating "it" with "the animal", but it is misleading. Because of value vectors and subsequent layers mixing the data, high attention weight does not definitively prove causal importance.
\end{enumerate}


% ─── Exercises ───
\section{Exercises}
\label{ch:09:sec:exercises}

\subsection*{Warmup}
\begin{enumerate}
  \item If $\vQ \in \R^{100 \times 64}$, how many operations (multiply-accumulates) are required to compute $\vQ\vK^\T$?
  \item In a causal mask, why do we use $-\infty$ instead of 0?
\end{enumerate}

\subsection*{Derivation}
\begin{enumerate}
  \item Prove mathematically why dividing by $\sqrt{d_k}$ ensures the variance of the dot product $\mathbf{q} \cdot \mathbf{k}$ is 1, given $q_i \sim \mathcal{N}(0, 1)$ and $k_i \sim \mathcal{N}(0, 1)$.
\end{enumerate}

\subsection*{Implementation}
\begin{enumerate}
  \item \textbf{[Prepares for CS336 A1]} Implement a PyTorch function \texttt{causal\_attention(q, k, v)} using raw tensor operations and \texttt{torch.einsum}. Verify that the attention weights in the upper triangle are exactly zero after the softmax.
\end{enumerate}

\subsection*{Open-Ended}
\begin{enumerate}
  \item Since causal attention restricts a token to only look backwards, does this mean token $x_{10}$ can never be influenced by token $x_{11}$ during backpropagation training? Think carefully about the flow of gradients.
\end{enumerate}

% ─── Further Reading ───
\section{Further Reading}
\label{ch:09:sec:further}

\begin{itemize}
  \item \textcite{vaswani2017}: Attention Is All You Need. The foundational paper that defined scaled dot-product attention and MHA.
  \item \textcite{bahdanau2014}: Neural Machine Translation by Jointly Learning to Align and Translate. The pre-Transformer origin of the attention mechanism in RNNs.
\end{itemize}
